\section{The remaining finite local laws on the original labels}
\label{sec:finite-local-completion-v43}

\status{The first part proves the full original root-injectivity
collision estimate, with its constant $23$, for every finite field
of cardinality at least $2^{12}$ and every degree.  The second part
works with the original finite projectors and normalized polynomial
labels.  Every use of Green's published cuspidal theorem is stated
explicitly.  The finite local statements do not themselves construct
the adelic rows or the arithmetic cylinder counts.}

\subsection{The fifth local error term}

Let $K$ be a finite field of cardinality $q$, and let $L$ be an
algebraically closed extension of $K$.  Write $\mathcal S_N$ for the
monic squarefree polynomials $F\in K[X]$ of degree $N$ with
$F(0)\ne0$, and put
\[
 w(F)=\prod_{P\mid F}(q^{\deg P}-1)^{-1},\qquad
 \mu_N(A)=\sum_{F\in\mathcal S_N\cap A}w(F).
 \tag{FLC.1}
\]
The factors $P$ are monic irreducibles.  This is exactly the
source's \texttt{sqWeight}, with the coefficient-vector
parametrization \texttt{monicOf} suppressed.  RRT5 proves
$\mu_N(\{F:H\mid F\})\le w(H)$ for monic squarefree $H$ with
nonzero constant term and $\deg H\le N$.  The same event is empty
if $\deg H>N$ or $H$ has a repeated irreducible factor.
For a prime $\ell$ and a nonzero finite-order element $z$, write
$z_{\ell'}$ for its prime-to-$\ell$ part.  Thus in the cyclic group
of order $\ell^a m$, with $(\ell,m)=1$, this is the projection
onto its subgroup of order $m$.  It commutes with every power map.
The original exceptional set is
\[
 \mathcal B(q,N)=\{q\}\cup
 \{\ell\text{ prime}:\exists\,1\le d\le N,
          q^d<\ell^{4v_\ell(q^d-1)}\}.
 \tag{FLC.2}
\]
Its finite-range implementation is identical: a prime satisfying
the displayed strict inequality divides $q^d-1$ and is at most
$q^N$.  No primality assumption on $q$ is needed in this subsection.

% BEGIN FLC1
\paragraph{FLC1. Frobenius orbits and the exact root-fiber bound.}
\label{claim:FLC1}
Every monic irreducible $P\in K[X]$ of degree $d$ with $P(0)\ne0$
has exactly $d$ distinct roots in $L$.  These roots form one orbit
under $z\mapsto z^q$, and each satisfies $z^{q^d-1}=1$.
For any integer $n\ge1$, prime $\ell$, and $a\in L^\times$,
\[
 \#\{z\in L:z^n=1,\ z_{\ell'}=a\}
                      \le\ell^{v_\ell(n)}.
 \tag{FLC.3}
\]
\begin{proof}
A root $z$ generates a field of $q^d$ elements over $K$.
The Frobenius orbit polynomial is fixed by Frobenius, hence has
coefficients in $K$, so its length is at least the degree of the
minimal polynomial.  Frobenius to the $d$th power is the identity
on this field.  Thus the orbit length is exactly $d$ and exhausts
the roots.  Finite fields are perfect: the $p$th-power map is an
injection and therefore a surjection on a finite field.  This also
proves separability.  Its nonzero elements have order dividing
$q^d-1$ by Lagrange's theorem.

If the set in (FLC.3) is nonempty, choose $z_0$ in it.
For every other member, $z/z_0$ has $\ell$-power order because the
prime-to-$\ell$ projections agree; its order also divides $n$.
It is therefore a root of $X^{\ell^{v_\ell(n)}}-1$.
The injective map $z\mapsto z/z_0$ and the polynomial root bound
prove (FLC.3), including the case where $\ell$ is the characteristic.
\end{proof}
% END FLC1

% BEGIN FLC2
\paragraph{FLC2. The full same-factor count.}
\label{claim:FLC2}
Let $d\ge2$, let $\ell$ be prime, and assume
$\ell^{4v_\ell(q^d-1)}\le q^d$.
The number of monic irreducible degree-$d$ polynomials for which
$z\mapsto z_{\ell'}$ is not injective on their roots is at most
$q^{3d/4}$.  This is the exact original \texttt{SameFactorCount}.
\begin{proof}
Put $a=\ell^{v_\ell(q^d-1)}\le q^{d/4}$.
An irreducible of degree at least two has nonzero constant term.
For a bad irreducible choose a root $z$ and a distinct conjugate
$z^{q^j}$ with equal projections.  The Frobenius orbit of
$z_{\ell'}$ has a proper length $e\mid d$, so $1\le e\le d/2$.
Consequently
\[
 (z_{\ell'})^{q^e-1}=1,\qquad z^{a(q^e-1)}=1.
 \tag{FLC.4}
\]
The second equation holds since the $\ell$-primary part of $z$
has order dividing $a$.  All $d$ roots of this irreducible
satisfy the same equation.  Different monic irreducibles have
disjoint root sets.  The union, over all proper positive divisors
$e$ of $d$, of the roots in (FLC.4) has cardinality at most
\[
 \sum_{\substack{e\mid d\\1\le e<d}} a(q^e-1)
                   \le d\,q^{d/4}q^{d/2}.
\]
Divide by $d$ to obtain the assertion.  The argument counts all
bad factors, without selecting a distinguished root of each one.
\end{proof}
% END FLC2

% BEGIN FLC3
\paragraph{FLC3. The full cross-factor count.}
\label{claim:FLC3}
Fix any irreducible $P'\in K[X]$ with $P'(0)\ne0$.
Let $d\ge1$ and $\ell$ be prime with
$\ell^{4v_\ell(q^d-1)}\le q^d$.
At most $q^{d/4}$ monic irreducible degree-$d$ polynomials $P$
with $P(0)\ne0$ have roots $z\in P^{-1}(0)$ and
$y\in(P')^{-1}(0)$ satisfying $z_{\ell'}=y_{\ell'}$.
This is the exact original \texttt{CrossFactorCount}.
\begin{proof}
Fix a root $y_0$ of $P'$.  By FLC1, for each such pair $(z,y)$
some common Frobenius power carries $y$ to $y_0$ and carries $z$
to another root of the same $P$.  This root lies in
\[
 T=\{x:x^{q^d-1}=1,\ x_{\ell'}=(y_0)_{\ell'}\}.
\]
Each relevant $P$ has a root in $T$, and roots belonging to
different monic irreducibles cannot coincide.  Therefore their
number is at most $|T|$.  FLC1 gives
$|T|\le\ell^{v_\ell(q^d-1)}\le q^{d/4}$.
No factor $\deg P'$ is lost: the common Frobenius transport is
what permits the single fixed fiber $T$.
\end{proof}
% END FLC3

% BEGIN FLC4
\paragraph{FLC4. Weighted coverage of both collision mechanisms.}
\label{claim:FLC4}
Let $A_N$ be the event that, for some prime
$\ell\notin\mathcal B(q,N)$, two distinct roots of $F$ have equal
prime-to-$\ell$ projections.  Then
\begin{align}
 \mu_N(A_N)&\le A+B,\notag\\
 A&=\sum_{d=2}^{N}\omega(q^d-1)
                \frac{q^{3d/4}}{q^d-1},\notag\\
 B&=\sum_{d=1}^{N}\sum_{e=1}^{d}
       \bigl(\omega(q^d-1)+\omega(q^e-1)\bigr)
                \frac{q^{d/4}}{q^d-1}.
 \tag{FLC.5}
\end{align}
Here $\omega(n)$ is the number of distinct prime factors of $n$.
\begin{proof}
Factor $F$ into its distinct monic irreducibles.  If both roots
belong to one factor of degree $d$, then $d\ge2$.  Also
$\ell\mid q^d-1$, since otherwise the projection is the identity
on all of its roots.  Being outside $\mathcal B(q,N)$ gives
the hypothesis of FLC2 for this $d$.  For each bad factor $P$,
RRT5 bounds the total weight of its multiples in $\mathcal S_N$
by $(q^d-1)^{-1}$.  Union bounds over factors and primes give $A$.

Otherwise the roots belong to distinct factors $P,P'$ of
respective degrees $d\ge e\ge1$.  Their collision forces
$\ell\mid(q^d-1)(q^e-1)$, since otherwise both projections are
identities.  Again the hypothesis of FLC3 holds for the larger
degree $d$.  For fixed $P'$ and $\ell$, FLC3 gives at most
$q^{d/4}$ possibilities for $P$.  The divisor bound assigns
weight at most $((q^d-1)(q^e-1))^{-1}$ to each distinct pair.
The number $I_e$ of monic irreducibles of degree $e$ with
nonzero constant term satisfies
$eI_e\le q^e-1$, by the disjoint-root count in FLC1; in
particular $I_e\le q^e-1$.  Summing first over $P'$, and then
over the at most $\omega(q^d-1)+\omega(q^e-1)$ possible primes,
gives the corresponding term of $B$.

Pairs with $P=P'$ or with $d+e>N$ contribute the empty event.
They can be omitted before applying the squarefree divisor bound.
Allowing them when bounding the number of pairs only increases
its nonnegative upper bound.  Thus no divisor estimate for a
nonsquarefree polynomial is being assumed.
\end{proof}
% END FLC4

% BEGIN FLC5
\paragraph{FLC5. The original constant $23$ and all original quantifiers.}
\label{claim:FLC5}
For every finite field $K$ with $q=|K|\ge2^{12}$, every
algebraically closed extension $L/K$, and every $N\ge0$,
\[
 \sum_{\substack{F\in\mathcal S_N\\
       \exists\ell\notin\mathcal B(q,N)\text{ prime}:\\
       z\mapsto z_{\ell'}\text{ not injective on roots of }F}}
       w(F)\le\frac{23\log q}{\sqrt q}.
 \tag{FLC.6}
\]
This proves the full original \texttt{LocalEstimateII}.
\begin{proof}
For $d\ge1$, the product of the distinct prime factors of
$q^d-1$ is at least $2^{\omega(q^d-1)}$ and divides $q^d-1$.
Hence
\[
 \omega(q^d-1)\le\frac{d\log q}{\log2}
                         \le\frac32d\log q.
\]
The numerical inequality $\log2>2/3$ follows, for example,
by integrating the strictly convex function $1/t$ on $[1,2]$:
its integral strictly exceeds its midpoint value $2/3$.
Put $x=q^{-1/4}\le1/8$.  Since $q^d-1\ge q^d/2$,
FLC4 gives
\[
 A\le3\log q\sum_{d=2}^N d x^d,\qquad
 B\le6\log q\sum_{d=1}^N d^2(x^3)^d.
 \tag{FLC.7}
\]
For $0\le x\le1/8$ and $0\le y\le1/512$, the elementary
inequalities $d\le2^d$ and $d^2\le4^d$ give
\[
 \sum_{d=2}^N d x^d\le\frac{4x^2}{1-2x}
                          \le\frac{16}{3}x^2,\qquad
 \sum_{d=1}^N d^2y^d\le\frac{4y}{1-4y}\le5y.
\]
Taking $y=x^3$ yields
\[
 A+B\le(16x^2+30x^3)\log q
       \le\frac{79}{4}x^2\log q
       \le23x^2\log q.
\]
Since $x^2=q^{-1/2}$, this is (FLC.6).  Empty degree ranges
cover $N=0$ and $N=1$ without an extra lower-rank hypothesis.
The definition of $\mathcal B$ and the canonical projection in
(FLC.2) are the original definitions, so this is the original
endpoint on the original polynomial space, not a substitute
collision event.
\end{proof}
% END FLC5

\subsection{Finite representation calculations with the original characters}

For the rest of the chapter, $q$ is prime and
$G_n=\operatorname{GL}_n(\mathbb F_q)$.  All representations are
finite dimensional over $\mathbb C$.  Fix a nontrivial additive
character $\psi$ of $\mathbb F_q$, and write
\[
 U_n=\{\text{upper unitriangular matrices}\},\qquad
 \lambda_n(u)=\psi\left(\sum_{a=0}^{n-2}u_{a,a+1}\right),\qquad
 \operatorname{Wh}(V)=V^{U_n,\lambda_n}.
 \tag{FLC.8}
\]
The superscript means the subspace on which $u$ acts as
$\lambda_n(u)$.  Its projector is exactly the $e_U$ of RRT1.
The rank-zero group and its character are trivial.  In particular
its Whittaker functor is the identity, with no exceptional sign.

% BEGIN FLC6
\paragraph{FLC6. The finite induction identities used below.}
\label{claim:FLC6}
Let $G$ be finite and let $H,K$ be subgroups.  For complex
representations $A$ of $H$ and $B$ of $K$,
\[
 \operatorname{Hom}_G(\operatorname{Ind}_H^G A,
                          \operatorname{Ind}_K^G B)
 \simeq\bigoplus_{x\in H\backslash G/K}
 \operatorname{Hom}_{H\cap xKx^{-1}}(A,{}^xB).
 \tag{FLC.9}
\]
This includes Frobenius reciprocity when $H=G$.
All representations are semisimple, and for a character
$\lambda:H\to\mathbb C^\times$,
\[
 \dim V^{H,\lambda}
       =|H|^{-1}\sum_{h\in H}\lambda(h)^{-1}\chi_V(h).
 \tag{FLC.10}
\]
\begin{proof}
A representation admits an invariant positive Hermitian form:
average any such form over $G$.  Orthogonal complements then
prove semisimplicity.  Realize an induced module as functions
$f:G\to A$ with $f(hg)=h f(g)$, with the right translation
action.  On restriction to $K$, separate these functions
according to their supports in the double cosets $HxK$.
Evaluation at $x$ identifies each summand with induction from
$K\cap x^{-1}Hx$ of the transported representation.  A map
from an induced module is uniquely determined by its values
on functions supported in $H$; conversely a map commuting with
$H$ extends to every coset by translation.  This proves
reciprocity, and applying it to the displayed support
decomposition proves (FLC.9), after replacing representatives
by their inverses if necessary.

The operator $|H|^{-1}\sum_h\lambda(h)^{-1}h$ is an idempotent
with image $V^{H,\lambda}$: multiplying by any $h_0\in H$
multiplies it by $\lambda(h_0)$, and it is the identity on that
subspace.  Taking its trace proves (FLC.10).
\end{proof}
% END FLC6

% BEGIN FLC7
\paragraph{FLC7. The exact Gelfand--Graev norm.}
\label{claim:FLC7}
For $n\ge1$, put $\Gamma_n=\operatorname{Ind}_{U_n}^{G_n}\lambda_n$.
Then
\[
 \dim\operatorname{Wh}(\Gamma_n)
       =\dim\operatorname{End}_{G_n}(\Gamma_n)
       =q^n-q^{n-1}.
 \tag{FLC.11}
\]
\begin{proof}
Gaussian elimination gives the Bruhat decomposition
$G_n=\coprod_{w\in S_n}U_n T_n w U_n$, and the diagonal factor
$t\in T_n$ is unique in the double quotient by $U_n$.
In detail, the pivot positions determine $w$; elimination using
unitriangular row and column operations then leaves the nonzero
pivot values $t$, which those operations cannot change.
Apply FLC6 with $H=K=U_n$.  The contribution of $tw$ is $1$
precisely when the two characters agree on
$U_n\cap twU_n(tw)^{-1}$, and is $0$ otherwise.

This intersection is generated by the elementary matrices
$1+aE_{ij}$ with $i<j$ and $w^{-1}(i)<w^{-1}(j)$.
The character is nontrivial on such a root subgroup exactly
when its two indices are adjacent.  Equality of the two
characters therefore requires
\[
 \begin{split}
 w(j)<w(j+1)&\ \Longrightarrow\ w(j+1)=w(j)+1,\\
 w^{-1}(i)<w^{-1}(i+1)&\ \Longrightarrow\
                  w^{-1}(i+1)=w^{-1}(i)+1.
 \end{split}
 \tag{FLC.12}
\]
The permutations satisfying these conditions are precisely
concatenations of increasing consecutive intervals, with the
intervals themselves arranged in decreasing order.
Here is an elementary verification.  Cut $w(1),\ldots,w(n)$
at each descent.  The first condition makes every resulting
run a consecutive interval.  If two such intervals occurred
in increasing value order, choose two successive value
intervals whose positional order increases; their boundary
values are consecutive, but occupy nonadjacent positions.
The second condition forbids this.  Thus the intervals occur
in decreasing order.  The converse follows immediately by
checking roots within and between intervals.

If there are $r$ intervals, character compatibility further
requires $t_i=t_{i+1}$ within each interval and imposes no
condition between intervals.  Indeed equality
$\psi(ca)=\psi(a)$ for every $a\in\mathbb F_q$ implies $c=1$:
otherwise $(c-1)a$ runs through the field, contradicting
nontriviality of $\psi$.  Hence this permutation contributes
$(q-1)^r$ compatible diagonal matrices.  Choosing the $r-1$
cuts among $n-1$ positions gives
\[
 \sum_{r=1}^n\binom{n-1}{r-1}(q-1)^r
                         =(q-1)q^{n-1}.
\]
Frobenius reciprocity identifies the endomorphism dimension
with $\dim\operatorname{Wh}(\Gamma_n)$, proving (FLC.11).
\end{proof}
% END FLC7

% BEGIN FLC8
\paragraph{FLC8. The generic mirabolic constituent and a sharp lower bound.}
\label{claim:FLC8}
Let
\[
 P_n=\left\{\begin{pmatrix}h&v\\0&1\end{pmatrix}:
              h\in G_{n-1},\ v\in\mathbb F_q^{n-1}\right\}
               \subset G_n.
\]
Then $R_n=\operatorname{Ind}_{U_n}^{P_n}\lambda_n$ is
irreducible.  Its dimension is
$\prod_{a=1}^{n-1}(q^a-1)$, and
$R_n|_{G_{n-1}}\simeq\Gamma_{n-1}$.
For every representation $W$ of $P_{k+1}$, with $k\ge1$,
\[
 \dim W^{U_k,\lambda_k}\ge
          (q^k-q^{k-1})\dim W^{U_{k+1},\lambda_{k+1}}.
 \tag{FLC.13}
\]
\begin{proof}
Write $V_n$ for the last-column normal subgroup of $P_n$.
Its characters are $v\mapsto\psi(av)$, with
$a\in(\mathbb F_q^{n-1})^\vee$; nontriviality of $\psi$
makes this parametrization injective and hence bijective.
The group $G_{n-1}$ has two orbits on these characters,
namely zero and nonzero.  The stabilizer of
$v\mapsto\psi(v_{n-2})$ is precisely $P_{n-1}$.

For completeness, the representation theory needed for this
observation is elementary.  The commuting finite-order
operators from $V_n$ are simultaneously diagonalizable.
Their eigenspaces are permuted by $G_{n-1}$.  An irreducible
$P_n$-module has a single orbit of nonzero eigenspaces.
Its chosen eigenspace is irreducible for the stabilizer;
otherwise translates of a proper invariant subspace would
produce a proper $P_n$-submodule.  Conversely, inducing an
irreducible stabilizer module, extended by the chosen scalar
character of $V_n$, is irreducible because any submodule is
the direct sum of its eigenspace intersections.  This proves
the orbit construction without invoking a classification.

Now $R_1$ is the trivial one-dimensional module.  Induction
in stages identifies $R_n$ with the nonzero-orbit induction
of $R_{n-1}$, since
$U_n=U_{n-1}\ltimes V_n$ and the last-column character is the
one just chosen.  The orbit construction proves irreducibility
by induction.  Counting the successive nonzero orbits gives
its stated dimension.  Moreover $P_n=G_{n-1}U_n$ and
$G_{n-1}\cap U_n=U_{n-1}$, so the one-double-coset restriction
in FLC6 gives $R_n|_{G_{n-1}}\simeq\Gamma_{n-1}$.

By reciprocity and irreducibility, $R_{k+1}$ is the unique
irreducible $P_{k+1}$-module with a nonzero
$(U_{k+1},\lambda_{k+1})$ space, and that space has dimension
one.  Thus a semisimple $W$ contains exactly
$m=\dim W^{U_{k+1},\lambda_{k+1}}$ copies of $R_{k+1}$.
Upon restriction to $G_k$, these contribute $m$ copies of
$\Gamma_k$.  All other summands contribute nonnegative
Whittaker dimensions.  FLC7 proves (FLC.13).
\end{proof}
% END FLC8

% BEGIN FLC9
\paragraph{FLC9. A finite character-exchange identity.}
\label{claim:FLC9}
Let $H_1,H_2\le H\le G$ be finite groups with
$|H_1|=|H_2|$, and let $\alpha_i:H_i\to\mathbb C^\times$
agree on $H_1\cap H_2$.
Suppose that, for $i=1,2$ and every $h\in H\setminus H_i$,
there are $a,a'\in H_i$ such that
$a'=h^{-1}ah$ and $\alpha_i(a')\ne\alpha_i(a)$.
Then for every class function $\chi$ on $G$,
\[
 |H_1|^{-1}\sum_{a\in H_1}\alpha_1(a)^{-1}\chi(a)
       =|H_2|^{-1}\sum_{b\in H_2}\alpha_2(b)^{-1}\chi(b).
 \tag{FLC.14}
\]
\begin{proof}
For $h\notin H_i$, the sum
$F(h)=\sum_{c\in H_i}\alpha_i(c)^{-1}\chi(ch)$ vanishes.
Indeed replace $c$ by $(a')^{-1}ca$ and use
$a'=h^{-1}ah$ and conjugation invariance; this gives
$F(h)=\alpha_i(a')\alpha_i(a)^{-1}F(h)$ with a scalar
different from one.
Now calculate
\[
 S=\sum_{a\in H_1}\sum_{b\in H_2}
        \alpha_1(a)^{-1}\alpha_2(b)^{-1}\chi(ab).
\]
First fix $b$.  Its contribution vanishes unless
$b\in H_1\cap H_2$.  In that case translation inside $H_1$
and agreement of the characters give the unshifted sum over
$H_1$.  Thus $S$ is $|H_1\cap H_2|$ times that sum.
Reversing the order and using $\chi(ab)=\chi(ba)$ gives the
corresponding expression for $H_2$.  Cancel the positive
intersection cardinality, then divide by the common group
order.  This proves the identity, with its exact normalized
averages.
\end{proof}
% END FLC9

% BEGIN FLC10
\paragraph{FLC10. The literal $e_J$ exchanges to the column-vanishing projector.}
\label{claim:FLC10}
Let $N=k+2s+1$, $k\ge1$, $s\ge0$, and use the original
$e_J=e_J(q,N,k,s,\psi)$ from RRT1.  Define
\[
 J'=\{u\in U_N:u_{a,k}=0\ (0\le a<k)\},\qquad
 e_{J'}=|J'|^{-1}\sum_{u\in J'}\lambda_N(u)^{-1}[u].
\]
Then $\operatorname{tr}(e_J\mid V)
=\operatorname{tr}(e_{J'}\mid V)$ for every representation of
$G_N$.  In fact the normalized averages agree for every class
function.  This is the full original pattern-exchange endpoint.
\begin{proof}
Use two ordered chains
\[
 W=(s,s+1,\ldots,s+k-1),\qquad
 C=(0,\ldots,s-1,s+k,s+k+1,\ldots,N-1),
\]
with entries $W_a$ and $C_b$ numbered from zero.  For
$0\le i\le s$, order the basis as
\[
 C_0,\ldots,C_{i-1},W_0,\ldots,W_{k-1},C_i,C_{i+1},\ldots,C_{2s}.
\]
Let $H_i$ consist of the matrices unitriangular in this order
whose $k$ entries from $W$ into column $C_i$ vanish.
Its permitted entries are all forward entries within $W$ and
within $C$, all $C_b\to W$ with $b<i$, and all $W\to C_b$
with $b>i$.  This is precisely the source pattern
\texttt{phI k s i}.  It is closed under multiplication: the
forbidden entry has no two-step path through permitted entries.
It has $q^{N(N-1)/2-k}$ elements.  Let $\alpha_i$ sum the
entries along the adjacent edges of the two fixed chains and
apply $\psi$.  No character edge has a two-step path through
permitted entries, so this sum is additive under matrix
multiplication.  These are exactly \texttt{eeI} and its
pattern character.

For $1\le i\le s$, put $p=C_{i-1}$ and $t=C_i$.
The patterns of $H_i$ and $H_{i-1}$ differ only as follows:
\[
 \begin{array}{c|c|c}
       &H_i\setminus H_{i-1}\text{ entries}
       &H_{i-1}\setminus H_i\text{ entries}\\ \hline
       &p\longrightarrow W&W\longrightarrow t.
 \end{array}
 \tag{FLC.15}
\]
The union pattern is a group $H$: insert $p,W,t$ in this
order and allow both displayed families.  The entry $p\to t$
is a character edge.  The two characters agree on the
intersection.

If $h\in H\setminus H_i$, some $h_{w,t}\ne0$ with $w\in W$.
Take $a=1+uE_{p,w}$.  Then $a\in H_i$, its character is one,
and
\[
 h^{-1}ah=1+u(h^{-1}e_p)(e_w^{\mathsf t}h)\in H_i,
 \qquad \alpha_i(h^{-1}ah)=\psi(u h_{w,t}).
\]
To check both assertions, $h^{-1}e_p$ is supported on the
$C$ entries up to $p$, whereas $e_w^{\mathsf t}h$ is supported
on $W$ from $w$ onward and on $C$ from $t$ onward.  All their
products are permitted in $H_i$; the only character edge
among them is $p\to t$, and $(h^{-1})_{p,p}=1$.
Choose $u$ so that the displayed character is not one.
Likewise, if $h\notin H_{i-1}$, then $h^{-1}\notin H_{i-1}$
and some $(h^{-1})_{p,w}\ne0$.  Use $a=1+uE_{w,t}$.
Now its conjugate lies in $H_{i-1}$ and has character
$\psi(u(h^{-1})_{p,w})$, again different from one for some $u$.
The support check is the same with the row and column roles
interchanged.  FLC9 therefore equates the two normalized
character averages.

Iterating from $i=s$ to $i=0$ proves the claim.  At $i=s$ the
order is the original order and $H_s=J$, with exactly
$\phi_J$ from RRT1.  At $i=0$, reordering the basis as $W,C$
identifies $H_0$ with $J'$ and its character with
$\lambda_N|_{J'}$.  Conjugating by this permutation matrix
does not change a class function.  For $s=0$ the single
permutation identification already proves the result.
\end{proof}
% END FLC10

% BEGIN FLC11
\paragraph{FLC11. The exact lower comparison on every original series.}
\label{claim:FLC11}
For every representation $V$ of $G_N$, with $N=k+2s+1$ and
$k\ge1$,
\[
 \dim e_JV\ge(q^k-q^{k-1})\dim e_UV.
 \tag{FLC.16}
\]
Consequently, for every central series projector $E_F$ that
is a sum of primitive central projectors,
\[
           (1-q^{-1})w_{e_U}(F)\le w_{e_J}(F).
 \tag{FLC.17}
\]
In particular this proves the full original
\texttt{lowerWtComparison}, on its actual \texttt{eSerC}
projectors, for every label, without Green's theorem.
\begin{proof}
Set $m=k+1$.  Let $T_m\subset U_N$ have identity upper-left
$m\times m$ block, arbitrary upper-right block, and upper
unitriangular lower-right block.  Define its character by
\[
 \beta(t)=\psi\left(\sum_{a=m-1}^{N-2}t_{a,a+1}\right).
\]
For $m=N$ this group and character are trivial.
The group $P_m$, embedded in the upper-left block of $G_N$,
normalizes $T_m$ and preserves $\beta$: its last row is fixed,
so the cross-block coordinate $t_{m-1,m}$ is unchanged.
Thus $W=V^{T_m,\beta}$ is a representation of $P_m$.
The actual semidirect products give
\[
 V^{U_N,\lambda_N}=W^{U_m,\lambda_m},\qquad
 V^{J',\lambda_N|_{J'}}=W^{U_k,\lambda_k}.
\]
Both equalities concern the original subspaces: $U_N=U_mT_m$,
$J'=U_kT_m$, and the indicated characters multiply to the
original adjacent-entry sums.  FLC8 and FLC10 prove (FLC.16).

For (FLC.17), apply (FLC.16) separately to each irreducible
in $E_F$.  Multiply by its positive degree, sum, and use RRT3's
coefficient identity together with
$[1]e_J/[1]e_U=q^k$ from RRT1.  The ratio of multiplicity
factors is $(q^k-q^{k-1})/q^k=1-q^{-1}$.
The argument works for any such central projector, so it
also applies directly to the original normalized projectors
without renaming their labels.
\end{proof}
% END FLC11

\subsection{Squarefree series and the unchanged normalized labels}

For a standard block parabolic $P=MU\subset G_n$, write
$\operatorname{HC}_P(\rho)=\operatorname{Ind}_P^{G_n}\rho$,
where $U$ acts trivially.  This is ordinary finite induction;
there is no modulus character.  An irreducible representation
of $G_d$ is cuspidal when its invariants under every proper
parabolic unipotent radical vanish.  It is enough to require
this for the two-block parabolics: a longer flag radical
contains a conjugate two-block radical.  The source's
\texttt{res a (d-a) chi 1 1=0} says precisely that the
averaging idempotent onto these invariants has trace zero.
Thus this definition is the original \texttt{IsCusp}.

% BEGIN FLC12
\paragraph{FLC12. Cuspidal degree determines the Whittaker multiplicity.}
\label{claim:FLC12}
For every cuspidal representation $V$ of $G_d$,
\[
 \dim V=\left(\prod_{a=1}^{d-1}(q^a-1)\right)
                         \dim\operatorname{Wh}(V).
 \tag{FLC.18}
\]
In particular, a cuspidal of degree
$\prod_{a=1}^{d-1}(q^a-1)$ has Whittaker multiplicity one.
\begin{proof}
For $1\le m\le d$, use exactly the tail subgroup $T_m$ and
character $\beta_m$ from FLC11, with ambient rank $d$,
and put $W_m=V^{T_m,\beta_m}$.
Then $W_d=V$, $W_1=\operatorname{Wh}(V)$, and $W_m$ is a
$P_m$-module.  For $m\ge2$, let $V_m$ be the last-column
normal subgroup of $P_m$.
The zero-character eigenspace $W_m^{V_m}$ vanishes: its
vectors are fixed by the entire unipotent radical of the
parabolic with blocks $(m-1,d-m+1)$.
Indeed the roots from the first $m-1$ rows into column $m-1$
belong to $V_m$, and those into later columns belong to $T_m$
and have trivial $\beta_m$; together they are that radical.
Cuspidality kills its invariant subspace.

All nonzero $V_m$-character eigenspaces in $W_m$ have equal
dimension, because $G_{m-1}$ acts transitively on their
characters.  The eigenspace for the last-coordinate character
is exactly $W_{m-1}$: adjoining $V_m$ to $T_m$ gives
$T_{m-1}$, and adds the single adjacent-entry term to its
character.  There are $q^{m-1}-1$ nonzero characters.  Hence
$\dim W_m=(q^{m-1}-1)\dim W_{m-1}$.
Multiplying these identities proves (FLC.18), including
$d=1$, when the product is empty.
\end{proof}
% END FLC12

% BEGIN FLC13
\paragraph{FLC13. Cuspidal support and the squarefree induction.}
\label{claim:FLC13}
Let $(\rho_i)_{i=1}^r$ and $(\rho'_j)_{j=1}^{r'}$ be lists
of cuspidal irreducibles, of positive ranks whose sums equal
$n$.  Their Harish--Chandra inductions satisfy
\[
 \dim\operatorname{Hom}_{G_n}
       (\operatorname{HC}(\rho_1,\ldots,\rho_r),
        \operatorname{HC}(\rho'_1,\ldots,\rho'_{r'}))
   =\#\{\sigma:\{1,\ldots,r\}\xrightarrow{\sim}
                       \{1,\ldots,r'\}:\rho_i\simeq\rho'_{\sigma(i)}\}.
 \tag{FLC.19}
\]
An isomorphism here includes equality of ranks.
Every irreducible of $G_n$ occurs in such an induction, and
the unordered list of cuspidal factors is uniquely determined.
If the factors are pairwise nonisomorphic, their induction
is irreducible.
\begin{proof}
For block sizes $(d_i)$ and $(e_j)$, the double cosets of the
two block parabolics are indexed by matrices of nonnegative
integers $(a_{ij})$ with row sums $d_i$ and column sums $e_j$.
One sees this directly from two flags: $a_{ij}$ is the dimension
of the successive quotient of their intersections.
Choosing bases for these quotients gives a simultaneous
adapted basis, and hence a representative; two pairs with
the same dimensions are carried to one another by a linear
isomorphism.  This proves both existence and completeness of
the double-coset indexing.

Apply FLC6.  If row $i$ has two or more nonzero entries, the
intersection subgroup on the $i$th Levi block includes the
unipotent radical of the proper flag with block sizes the
nonzero $a_{ij}$.  This subgroup acts trivially on the other
inducing representation, because there it lies in the
parabolic radical; its matrices are the identity on every
other block of the first Levi.  An intertwiner must consequently factor
through the invariants of $\rho_i$ under this proper
radical: average over that subgroup to see this explicitly.
Those invariants vanish by cuspidality.  The same argument,
with the flags interchanged, applies to a split column.
Only matrices having one nonzero entry in each row and each
column survive.  They are exactly rank-matching bijections
$\sigma$, and their intertwining spaces are
$\bigotimes_i\operatorname{Hom}_{G_{d_i}}(\rho_i,\rho'_{\sigma(i)})$.
Schur's lemma makes their dimensions zero or one, proving
(FLC.19).  Schur's lemma itself follows because the kernel
and image of an intertwiner are invariant, and any complex
endomorphism has an eigenvalue.

For existence of a cuspidal support, start with an irreducible
$V$.  If it is not cuspidal, some proper two-block radical
has nonzero invariants.  Choose an irreducible constituent
of this Levi module.  Irreducibles of a product are tensor
products, as follows by decomposing under one factor and
using its commuting multiplicity spaces.  Reciprocity embeds
$V$ as a constituent of the corresponding induction.
Repeat in smaller positive ranks until all factors are
cuspidal.  Induction in stages, proved by composing the
function models, gives one induction from the resulting list.
If two such inductions share $V$, semisimplicity gives a
nonzero map between them, so (FLC.19) identifies their lists.
Finally a list of pairwise nonisomorphic factors has just
the identity self-matching; its induced module has
endomorphism dimension one and therefore is irreducible.
The empty list gives the trivial rank-zero representation.
\end{proof}
% END FLC13

% BEGIN FLC14
\paragraph{FLC14. Whittaker multiplicity is multiplicative in finite induction.}
\label{claim:FLC14}
For representations $V_i$ of $G_{d_i}$,
\[
 \dim\operatorname{Wh}(\operatorname{HC}(V_1,\ldots,V_r))
                    =\prod_{i=1}^r\dim\operatorname{Wh}(V_i).
 \tag{FLC.20}
\]
\begin{proof}
The double cosets $U_n\backslash G_n/P$ are represented by
shuffles of the ordered Levi blocks.  To verify this, Gaussian
elimination gives $G_n=\coprod_w B_n wP$, with $w$ preserving
the order within each block.  Since diagonal matrices can
be moved through $w$ into $P$, the same representatives give
$U_n\backslash G_n/P$.
Let $b_a$ be the block occupying output coordinate $a$ of
such a shuffle.  If $b_a<b_{a+1}$, the elementary root subgroup
$1+tE_{a,a+1}$ lies in the conjugate parabolic radical.
It acts trivially on the inducing module but nontrivially
through $\lambda_n$.  Its term in reciprocity is therefore
zero.  The only shuffle without such an increase has the
blocks in decreasing order, with each block contiguous.
For this shuffle, the intersection with $U_n$ has the
within-block groups $U_{d_i}$, with their original adjacent
characters; every remaining radical coordinate has trivial
character.  Its intertwining space is the tensor product
of the $\operatorname{Wh}(V_i)$.  FLC6 proves (FLC.20).
This also proves the formula for character traces and all
parenthesizations of \texttt{hcL}, by induction in stages.
\end{proof}
% END FLC14

% BEGIN FLC15
\paragraph{FLC15. Elliptic character values uniquely determine the prime label.}
\label{claim:FLC15}
Fix $q\ge5$ and the original \texttt{LabelDatum} $L_0$:
its injective multiplicative Teichmuller map is $\iota_0$,
and its degree-$d$ generator is $\gamma_d$, of order $q^d-1$.
For a root $x$ of a prime label of degree $d$, let
\[
 \theta_x(\gamma_d^a)=\iota_0(x)^a.
 \tag{FLC.21}
\]
If the same cuspidal character satisfies the original
\texttt{IsCuspOf} formula for two prime labels $P,P'$ of
degree $d$, then $P=P'$.
\begin{proof}
The expression (FLC.21) is independent of $a$ modulo
$q^d-1$ and is exactly the source's \texttt{thetaOf}.
Choose roots $x$ and $x'$ of the two labels.  On
$\mathbb F_{q^d}^{\times}$ the difference
\[
 f(y)=\sum_{i=0}^{d-1}\theta_x(y^{q^i})
              -\sum_{i=0}^{d-1}\theta_{x'}(y^{q^i})
\]
vanishes whenever $y$ has degree $d$ over $\mathbb F_q$.
Indeed multiplication by such a $y$ on
$\mathbb F_{q^d}$ is a matrix with irreducible characteristic
polynomial, and its elliptic character value is the same
under the two given formulas.  The common sign
$(-1)^{d-1}$ is nonzero and cancels.

Let $B$ be the number of nonzero elements of degree less than
$d$ in $\mathbb F_{q^d}$.  Then
\[
                         2dB<q^d-1.                 \tag{FLC.22}
\]
For $d=1$, $B=0$.  For $d=2$, $B=q-1$, so the assertion is
$4<q+1$.  For $d\ge3$ put $h=\lfloor d/2\rfloor$.
Every proper subfield has degree a proper divisor of $d$,
so $B<q^{h+1}/(q-1)$ by a geometric sum.
The inequality $q^{\lceil d/2\rceil-1}\ge d$ follows for
$d=3,4$ from $q\ge5$, and then follows for $d+2$ by
multiplying by $q\ge5$.  Therefore
$2dB<2q^d/(q-1)\le q^d/2<q^d-1$, proving (FLC.22).

A nonzero function on a finite abelian group of order $M$
whose Fourier expansion uses at most $r$ characters has
support of cardinality at least $M/r$.  To prove this,
write $f=\sum c_\eta\eta$.  Orthogonality gives
$\sum_y|f(y)|^2=M\sum_\eta|c_\eta|^2$, whereas
$|f(y)|^2\le r\sum_\eta|c_\eta|^2$ by Cauchy--Schwarz.
Summing only over the support and canceling proves the claim.
Here $M=q^d-1$ and $r\le2d$, so (FLC.22) forces $f=0$.
Character orthogonality now identifies its two multisets of
characters.  Thus $\theta_{x'}=\theta_{x^{q^i}}$ for some $i$.
Evaluating at $\gamma_d$ and using injectivity of $\iota_0$
on nonzero elements gives $x'=x^{q^i}$.  Their monic minimal
polynomials are equal, proving $P=P'$.
\end{proof}
% END FLC15

\paragraph{The one published finite-group input.}
We use Green's cuspidal theorem in the following precise form.
For every finite field $\mathbb F_q$ and $d\ge1$, the regular
characters $\theta:\mathbb F_{q^d}^{\times}\to\mathbb C^\times$,
modulo Frobenius, parametrize the cuspidal irreducibles of
$G_d$.  Regular means that the Frobenius orbit has length $d$.
Their degrees and elliptic values are
\[
 \dim\rho_\theta=\prod_{a=1}^{d-1}(q^a-1),\qquad
 \chi_{\rho_\theta}(g)=(-1)^{d-1}
                     \sum_{i=0}^{d-1}\theta(y^{q^i})
 \tag{FLC.23}
\]
when $g$ has irreducible characteristic polynomial with root
$y\in\mathbb F_{q^d}$.
The published input is
J.~A.~Green, \emph{The characters of the finite general linear
groups}, Trans. Amer. Math. Soc. \textbf{80} (1955), 402--447:
Theorems 12--13, Lemma 7.4, and Theorem 9.
\footnote{\url{https://www.ams.org/tran/1955-080-02/S0002-9947-1955-0072878-2/S0002-9947-1955-0072878-2.pdf}.}
Its exact dictionary is as follows.  Writing $\theta=\theta_0^a$
for a primitive character, the regular orbit is Green's
degree-$d$ simplex $g=\{a,aq,\ldots,aq^{d-1}\}$ modulo $q^d-1$.
Theorem 12 gives
\[
 \rho_\theta=(g^{(1)})=(-1)^{d-1}I_d^a[(1)].
\]
The paper's $\circ$-product is our finite Harish--Chandra induction.
Its Section 7 Schur-function homomorphism preserves simplex
support, with $s_{(1)}^v=\sum_{\lambda\vdash v}f^\lambda s_\lambda$
and $f^\lambda>0$.  Thus $(g^{(1)})$ cannot occur in proper
induction; every other symbol does, by factoring distinct
supports or this identity.  Reciprocity and Theorem 13 identify
exactly the cuspidals.  Definition 3.1 and (12), page 417,
with Theorem 9, give
\[
 I_d^a[(1)](g_y)=T_{d,1}(a;y)
                    =\sum_{i=0}^{d-1}\theta_0^{a q^i}(y).
\]
Lemma 7.4 gives
$I_d^a[(1)](1)=\prod_{j=1}^{d-1}(1-q^j)$.
Both formulas, with the displayed sign, are (FLC.23).
This applies Green's classification; it does not reprove it.

% BEGIN FLC16
\paragraph{FLC16. Green's theorem gives the literal \texttt{Green} record.}
\label{claim:FLC16}
For every prime $q$, every original \texttt{LabelDatum} $L_0$,
and every $d\ge1$, there is a bijection
\[
 \{P\text{ monic irreducible}:\deg P=d,\ P(0)\ne0\}
       \longleftrightarrow
       \{\text{cuspidal irreducible characters of }G_d\}
 \tag{FLC.24}
\]
whose character $\tau_P$ satisfies \texttt{IsCuspOf} for that
same $L_0$ and the degree formula in (FLC.23).
This supplies every field of the original \texttt{Green}
predicate using the stated published theorem.
\begin{proof}
Work in the unique degree-$d$ subfield of the fixed algebraic
closure.  The restriction of $\iota_0$ to its multiplicative
group is a bijection onto the complex roots of unity of
order dividing $q^d-1$: it is injective and both finite sets
have $q^d-1$ elements.  Since $\gamma_d$ is a generator,
(FLC.21) therefore identifies elements $x$ of this group
with all its complex characters.  It intertwines Frobenius
on $x$ with Frobenius on characters.  Characters with orbit
length $d$ correspond exactly to elements of degree $d$,
and their Frobenius orbits correspond exactly to the monic
irreducible polynomials of degree $d$ with nonzero constant
term.  Apply (FLC.23) to these characters.

Choosing any root of $P$ gives the same cuspidal character,
because the published parameter is a Frobenius orbit.
Every root of $P$ is a Frobenius conjugate, so the elliptic
formula holds for each root, exactly as required by
\texttt{IsCuspOf}.  Roots of an irreducible characteristic
polynomial of a rank-$d$ matrix also lie in this same finite
subfield, so its quantified elliptic arguments are unchanged.
Formula (FLC.21) is the literal \texttt{thetaOf}; independence
of the chosen exponent was proved in FLC15's first sentence
and does not require $q\ge5$.
The published injectivity and surjectivity on Frobenius
orbits give the remaining two fields of \texttt{Green}.
The supplied $\gamma_d$, $\iota_0$, and their norm compatibility
are retained.  No replacement compatible generator system
or new embedding is selected.
\end{proof}
% END FLC16

% BEGIN FLC17
\paragraph{FLC17. The exact squarefree $e_U$ law in character labels.}
\label{claim:FLC17}
Let $q\ge5$ be prime and use any original \texttt{LabelDatum}.
For a monic squarefree degree-$N$ polynomial $F$ with
$F(0)\ne0$, the original \texttt{eSerC} series consists of
one irreducible $\tau_F$, and
\[
 \dim\operatorname{Wh}(\tau_F)=1,\qquad
 \dim\tau_F=
       \frac{\prod_{a=1}^{N}(q^a-1)}
                         {\prod_{P\mid F}(q^{\deg P}-1)},\qquad
 w_{e_U}(F)=\prod_{P\mid F}(q^{\deg P}-1)^{-1}.
 \tag{FLC.25}
\]
This is the full original squarefree Gelfand--Graev weight law.
\begin{proof}
FLC16 assigns a distinct cuspidal $\rho_P$ to every distinct
factor $P\mid F$.  FLC15 says that the source's
\texttt{IsCuspOf} relation has no other prime label for a
cuspidal.  Thus FLC13's support uniqueness identifies the
original \texttt{HasLabel} series of $F$ with the constituents
of $\operatorname{HC}((\rho_P)_{P\mid F})$.
The factors are pairwise nonisomorphic, so FLC13 makes this
induction irreducible, proving the singleton assertion.
Its Whittaker multiplicity is one by FLC12, FLC14, and the
published cuspidal degree supplied in FLC16.

For completeness, its degree uses only the index of the
actual block parabolic.  A parabolic with blocks of sizes
$d_1,\ldots,d_r$ has order
$q^{\sum_{i<j}d_id_j}\prod_i|G_{d_i}|$.
Counting ordered bases gives
$|G_n|=q^{n(n-1)/2}\prod_{a=1}^n(q^a-1)$.
Multiplying the parabolic index by
$\prod_i\prod_{a=1}^{d_i-1}(q^a-1)$ cancels the powers of
$q$ and gives the displayed degree.  Finally RRT3 gives
\[
 w_{e_U}(F)=\frac{\dim\tau_F}{|G_N|[1]e_U}
          =\frac{\dim\tau_F}{\prod_{a=1}^{N}(q^a-1)},
\]
since $\dim e_U\tau_F=1$ and $[1]e_U=q^{-N(N-1)/2}$.
This proves the exact normalization in (FLC.25).
For $N=0$, the empty support and all empty products equal one,
which is also the original rank-zero statement.
\end{proof}
% END FLC17

In the original datum $D$, the stored type-prime congruence is
\[
D.\mathrm{host}.g(k)\equiv1\pmod{D.\mathrm{host}.QS\,
D.\mathrm{host}.q\,D.\mathrm{host}.d}.
\]
The stored tower divisibility makes this modulus divisible by
$4$.  Thus each original type prime is $1$ modulo $4$ and,
being prime, is at least $5$.  The preceding character-label
hypotheses therefore apply at both of the original type primes.

% BEGIN FLC18
\paragraph{FLC18. The original normalized-label comparison.}
\label{claim:FLC18}
\mbox{}\par\noindent
Fix an actual \texttt{NormalizedTypesCompatibility} $T$ for
the original datum $D$ and its specified embedding $j_E$.
At either type prime $q=D.\mathrm{host}.g(k)$, let $O=T.O(k)$
be its original \texttt{LabelMap}.  Let $L_0$ use exactly
$D.\mathrm{ntypes}.\mathrm{teichC}(k)$ and $T.\zeta(k)$.
There is a degree-preserving permutation $\sigma$ of prime
polynomial labels such that, for every valid degree-$N$ label,
\[
 \widehat\sigma(F)=\prod_{P\mid F}\sigma(P)^{v_P(F)},\qquad
 E_F^{\mathrm{original}}=e\mathrm{SerC}_{L_0}(\widehat\sigma(F)).
 \tag{FLC.26}
\]
It preserves squarefreeness and the multiset of factor degrees.
Consequently the exact squarefree law (FLC.25) and lower
comparison (FLC.17) hold on the original labels themselves.
\begin{proof}
Fix $d\ge1$.  For each Green prime label $P$ choose a
representation $V_P$ of its cuspidal character from FLC16.
Send $P$ to $O.\mathrm{lab}(d,V_P)$.
The original \texttt{lab\_spec} makes this another prime
label of degree $d$.  This map is injective: equal $O$-labels
give isomorphic representations by the original
\texttt{lab\_inj}, hence equal Green characters and equal
$P$.  The finite set of degree-$d$ prime labels maps
injectively to itself, so this map is a permutation $\pi_d$.
Take $\sigma_d=\pi_d^{-1}$.

For any cuspidal representation $V$, FLC16 gives a $P$ with
its character.  Equal characters imply isomorphic complex
representations: averaging on $\operatorname{Hom}(V,V_P)$
computes its invariant dimension as their character pairing,
which is one by Schur's lemma and semisimplicity.
Thus \texttt{lab\_iso} gives
$O.\mathrm{lab}(d,V)=\pi_d(P)$ and
$\sigma_d(O.\mathrm{lab}(d,V))=P$.
This constructs the same degree-preserving prime-label
permutation required in the source comparison.
Its extension to products preserves degrees and
multiplicities; unique factorization gives
$\widehat\pi\widehat\sigma(F)=F$ and preservation of
squarefreeness.

The original \texttt{HasLabelO} says exactly that an
irreducible is a constituent of a block-parabolic induction
of cuspidals and that $F$ is the product of their $O$-labels.
Transporting the original block indices to consecutive
indices conjugates the groups and representations; finite
induction is carried by the same change of basis, so its
character is the source's \texttt{hcL}.  No induction
normalization changes, since the group is finite.
FLC13's unique cuspidal support and the preceding prime-label
identity therefore give, for each irreducible $V$,
\[
 \mathrm{HasLabelO}(O,V,F)
 \quad\Longleftrightarrow\quad
 \mathrm{HasLabel}(L_0,\chi_V,\widehat\sigma(F)).
\]
FLC15 ensures that the right-hand label relation is the
same one used in the source.  Summing the same primitive
central idempotents proves (FLC.26), using the supplied
literal equality $T.\mathrm{eSer\_def}$.
In particular, this is an equality of the original group
algebra elements, not just of their ranks or eigenvalues.

The squarefree weight in (FLC.25) depends only on the factor
degrees, which $\widehat\sigma$ preserves.  The lower law
holds for every series and transfers through (FLC.26).
These comparisons are made at the original polynomial $F$;
no root-collision predicate is transported by an arbitrary
permutation of labels.  Subsequent polynomial estimates
therefore still use the actual roots of $F$.
\end{proof}
% END FLC18

% BEGIN FLC19
\paragraph{FLC19. The retained local weight producer on the original row object.}
\label{claim:FLC19}
Fix the original datum $D$, the original
\texttt{NormalizedTypesCompatibility} $T$ of FLC18, and an
actual \texttt{AllRows} object $A$ with
\texttt{AllJoinA A}.  Retain the original
\texttt{TieAdmG} and \texttt{RangeW} hypotheses.
Write $f=D.R.f$, $m=D.R.m$, $N=D.R.r$, and let $N_0$ be the
absolute threshold in RRT9--RRT10.  If $N\ge N_0$, then the
actual assertion \texttt{WeightAllRows A} holds, with its
unchanged constant
\begin{align}
 W_g={}&\frac{3f}{q}+\frac{23f\log q}{\sqrt q}
       +2fm\left(\frac1{4\min(q_1,d_0)}
                     +\frac{1+\log N}{q-1}\right)\notag\\
     &+\frac N2(7/8)^f+\exp\{-f/(5\log N)\},
       \qquad q=D.\mathrm{host}.g(k).
 \tag{FLC.27}
\end{align}
Here $(q_1,d_0)=(D.\mathrm{host}.q,D.\mathrm{host}.d)$.
The conclusion contains all local nonnegativity assertions
and the weighted bound for every finite set of bad tuples,
on the original $A.eVec$ and original \texttt{goodAt}.
For the original selected eventual family, the additional
$N_0$ guard is absorbed by its actual rank growth, so the
original eventual retained-defect application has no
remaining finite local quantitative input beyond the stated
published theorem and the original supplied normalized-type
record.
\begin{proof}
The exact guards supplied by the two original records include
$q\ge2^{12}$, distinct odd primes $q_1,d_0$ with
$\max(q_1,d_0)\ge16$, and $16,q_1,d_0\mid q-1$.
The original \texttt{AllJoinA.level} gives, at each type
prime, one nontrivial additive character and the literal
$e_J$ at every place.  More generally the argument permits
the exact \texttt{TieVecW} alternatives $e_U$ or $e_J$.
RRT1--RRT2 give their nonnegative normalized probability
laws on the actual degree-$N$ labels, with weight zero on
invalid labels.  FLC18 supplies the squarefree $e_U$ formula
on those original labels, and FLC11 supplies the required
lower $e_J$ comparison.

Apply RRT11, replacing its last collision input by FLC5.
Its five bad-event mechanisms now all have proofs:
RRT6 handles nonsquarefree factors; FLC5 handles failure of
root injectivity; RRT16 handles each of the $2fm$ root-to-twist
collisions; RRT8 handles a common proper divisor degree;
and RRT10 handles absence of the long prime-degree witness.
RRT4 transfers the resulting bound to the actual local
projectors at a cost at most $f/q$.  This is exactly (FLC.27).
The original factor count $|D.\mathrm{nd}.\mathrm{PlG}(k)|=f$
and twist count $m$ give its coefficients.
Any finite set of bad tuples is contained in the full bad
event and has no larger weight.  Valid polynomials are
identified injectively by their coefficients; invalid ones
have zero projector.  These facts give precisely the two
clauses of \texttt{WeightAllRows}, even when different global
rows have the same local label.

The original consumer \path{ConcNWeight.locShare_le_Wg}
uses only the first three \texttt{admLq} clauses.
We have proved those clauses directly for the actual
projectors.  The fourth clause, the upper $e_J$ comparison,
is unused in this argument.

For the actual selected family, the original rank record
has $f=2\operatorname{ord}\le k\le r=N$, and its size
choice has $x^2\le f$.  Thus $N\ge N_0$ whenever
$x\ge\sqrt{N_0}$, uniformly in every later local choice.
This is the same legitimate rank comparison written in
RRT11.  It does not follow from the abstract upper/lower
bounds on $N$ in \texttt{SizesAt} alone.
\end{proof}
% END FLC19

\subsection{Exact scope of the finite completion}

FLC1--FLC5 close the full original \texttt{SameFactorCount},
\texttt{CrossFactorCount}, and \texttt{LocalEstimateII}.
FLC6--FLC11 close the actual pattern exchange and the full
lower local character comparison.  FLC12--FLC18 give the
squarefree series law and the degree-preserving normalized
label comparison, using the explicitly stated published
cuspidal theorem.  Every polynomial collision condition is
then tested on the original label, not on a relabeled root
set.  FLC19 supplies the actual retained local weight clause
for the selected eventual family with the original $W_g$.

This chapter retains the original supplied normalized-type
compatibility record; it does not construct that record from
bare abstract projectors.  It does not produce the original
adelic \texttt{HidaSetting}, the actual reduction and Haar
normalization in \texttt{EtOne}, the original
\texttt{RowsAdd} and \texttt{RowsGood} identities, the common
\texttt{AllRows} object, or any of the three sufficient
arithmetic cylinder counts in GQT5.  The full fourth
\texttt{admLq} clause, an upper comparison unused by this
consumer, is not claimed here.  The long-prime-degree term
still uses RRT10's eventual theorem, with its explicit
actual-rank justification.  These scope statements prevent
a local completion from being mistaken for completion of
\texttt{GLOBAL\_RETAINED\_DEFECT}, RH, or GRH.
